\section{Technical overview}\label{sec:overview}

This section explains the main ideas of Algorithm~1 and its analysis without tracking explicit precision constants, giving forward pointers to the formal statements in Sections~\ref{sec:prelim}--\ref{sec:lean}.

\subsection{The algorithm in five steps}\label{sec:overview-five-steps}

Below the fixed cutoff $n_*$ of~\eqref{eq:cutoff}, Algorithm~1 delegates directly to the independent inner multiplier $\mathsf M$. For $n\ge n_*$, set $b=\lfloor\log_2 n\rfloor$, $N_{\mathrm{dig}}=\lceil n/b\rceil$, and let $R=2^k$ be the least power of two with $R\ge 2N_{\mathrm{dig}}$. Factor $R=t_1\cdots t_d$ into balanced powers of two $\mathbf t=(t_1,\dots,t_d)$ with $d=\lfloor k^\eps\rfloor$, $t_d=r=2^\ell$, $t_i\in\{r/2,r\}$, and $A_a=\prod_{i<a}t_i$ ($1\le a\le d+1$). Algorithm~1 multiplies two $n$-bit integers on this single power-of-two grid in five steps.
\begin{enumerate}[label=\arabic*.,leftmargin=*,nosep,topsep=3pt]
  \item \emph{Complex digit packing} (Proposition~\ref{prop:root-packing}). Split the two inputs into radix-$2^b$ digits $x^{(0)}_j,x^{(1)}_j\in[0,2^b)$ ($0\le j<N_{\mathrm{dig}}$) and pack both into one complex vector $x_j=2^{-b}(x^{(0)}_j+i x^{(1)}_j)$ of length~$R$ (padded with zeros for $j\ge N_{\mathrm{dig}}$), reducing integer multiplication to two forward length-$R$ DFTs at working precision $P=\Theta(k)=\Theta(\log n)$ ($q=P+2$).
  \item \emph{Descending fractional-shift spreading} (Section~\ref{sec:bridge}). Pre-scale into $\Cg_q$ and spread axes $a=d,\dots,2$ in descending order on the same $t_1\times\cdots\times t_d$ grid by the kernel $g_a(s)=e^{-\pi i s/t_a}\sinc(s)e^{-\pi s^2/\alpha^2}$ ($\alpha=2\sqrt{\tw\ln 2/\pi}$, $\tw=\Theta(k)$) shifted by $\delta_a(\mathbf j)=(\sum_{i<a}j_i A_i)/A_a$ ($\mathbf j=(j_1,\dots,j_d)$, $j=\sum_{i=1}^d j_i A_i$, $0\le j_i<t_i$) and folded modulo~$t_a$ (Definition~\ref{def:carrier}, Lemma~\ref{lem:bridge-A}; exposures by Lemma~\ref{lem:bridge-D}), producing $U$ and absorbing Cooley--Tukey twiddles without Nyquist singularities.
  \item \emph{Direct complex tensor FFT} (Sections~\ref{sec:walsh}--\ref{sec:direct-fft}). Evaluate $Y\approx F_{\mathbf t}U$ on $\Cg_q$ directly over~$\C$, executing the chunk and bit-reversal layouts $L,(LB)^{-1}$ bit-exactly via three-layer selective Walsh permutations (Theorem~\ref{thm:positional-perm}) and the record-axis butterfly rounds with combined twiddles (Theorem~\ref{thm:direct-fft}).
  \item \emph{Ascending two-term cyclic deconvolution} (Lemmas~\ref{lem:bridge-B}--\ref{lem:bridge-C}). Invert the non-commuting two-alias operators $B_2,\dots,B_d$ in ascending order $a=2,\dots,d$ by solving $c_0 Z_f + c_1 Z_{(f-R/A_a)\bmod R}=Y_f$ on cycles of length $A_a$ (with $c_\nu=\hat g_a(u_a/t_a-\nu)$ for $f=\sum_{i=1}^d u_i R/A_{i+1}$ and $\|B_a^{-1}\|_\infty\le\frac47 t_a$ uniform in $A_a$), updating $Y\leftarrow Z$ after each axis, and rescale to obtain $\hat y\approx F_R x$.
  \item \emph{Conjugated squaring and carry propagation} (Proposition~\ref{prop:root-packing}). Form $\hat w=\overline{\Quant_{p+2}(\hat y\odot\hat y)}$, repeat Steps~2--4 on $\hat w$ to obtain $\hat z\approx F_R\hat w$, recover the exact integer convolution $x^{(0)}\circledast x^{(1)}$ by rounding $-2^{2b-1}R^2\,\mathrm{Im}\,\hat z_j$ via the two-call identity~\eqref{eq:two-call}, and propagate radix-$2^b$ carries.
\end{enumerate}

\subsection{Where the \texorpdfstring{$n\log n$}{n log n} barrier is broken, and what caps the saving}\label{sec:overview-barrier}

Because $R=\Theta(n/\log n)$ and working precision $P=\Theta(k)=\Theta(\log n)$ suffices to round the output coefficients to exact integers, the stored array has logical bit volume $V=\Theta(RP)=\Theta(n)$. On a multitape Turing machine, a standard one-dimensional FFT performs $k=\Theta(\log n)$ sequential passes over volume~$V$, costing $\Theta(n\log n)$ steps from tape movement alone.

Harvey and van der Hoeven~\cite{HvdH21} reach $O(n\log n)$ by folding a one-dimensional convolution of odd-prime-product length $S=s_1\cdots s_d$ into $d$ dimensions. A Chinese remainder map splits the length onto the pairwise coprime odd primes $s_1,\dots,s_d$, Gaussian resampling interpolates each prime axis onto a larger power-of-two length $t_i$, and multidimensional FFTs evaluate the product. In~\cite{HvdH21} the dimension~$d$ is a large constant. In~\cite{OAI26}, by contrast, $d=\lfloor k^\eps\rfloor$ grows with~$n$, and two fixed finite linear networks at $h=100$ ($m=h^3=10^6$; Fact~\ref{fact:oai-networks}) give a power saving across the $d$ axes.
\begin{enumerate}[label=(\roman*),nosep,leftmargin=*]
  \item A \emph{bit-shear network} of $W_b$ roles with $s_b < W_b m^{\tau}$ ($\tau<1$) interchanges two disjoint contiguous $u$-bit address fields in time $O(V u^{\tau})$ rather than $O(V u)$ (Lemma~\ref{lem:transfer12}).
  \item A \emph{complex phase network} of $W_c$ roles with $s_c < W_c m^{\sigma}$ ($\sigma<1$) executes the normalized butterfly $H_0=\frac12\bigl(\begin{smallmatrix}1&1\\1&-1\end{smallmatrix}\bigr)$ simultaneously on one selected bit~$\rho$ of each of $D\le d$ consecutive $K$-bit address chunks ($K=\lfloor d^c\rfloor$) in time $O(V d^{\lambda'})$ with $\lambda'<1$, rather than $O(V d)$ (Proposition~\ref{prop:prop14-param}).
\end{enumerate}
Since $\ell=\lceil k/d\rceil$ bit-levels across $d-1$ record axes then cost $O(V\ell d^{\lambda'})=O(V k d^{\lambda'-1})$, one would hope for an exponent saving $\kappa\approx\eps(1-\lambda')$ linear in the network deficit $1-\sigma$.

\paragraph{The chunk-layout cap.}
Why does~\cite{OAI26} instead obtain $\kappa=2^{-182}$? Before the simultaneous butterfly layer can act, the $m_0=\log_2(R/r)$ record-address bits, initially stored axis by axis, must be rearranged by a layout permutation~$L$ that interleaves $K$-bit chunks across axes. Executing~$L$ via $O(k/K)$ width-$K$ chunk interchanges costs $O\bigl(V\cdot (k/K)K^{\tau}\bigr)=O(V P K^{\tau-1})$, whose saving factor $K^{1-\tau}=\Theta(P^{\eps c(1-\tau)})$ requires $c$ to be large. Yet inside the butterfly layer recursion (stopping below $d^\beta$ chunks), moving $K$-bit chunks into aligned basis positions forces $\tau(1+c/\beta)<\lambda<1$, hence $c<\beta(1-\tau)/\tau$. Meanwhile, keeping the internal precision guard $\Delta=\lceil C_0 d^{C_1}\rceil$ sublinear in~$P$ forces $\eps<1/C_1$. Multiplying these bounds traps any chunk-exchange layout beneath a \emph{quadratic} ceiling:
\begin{equation}\label{eq:overview-quad-cap}
  \kappa < \eps c(1-\tau) < \frac{\beta(1-\tau)^2}{C_1\tau} < \frac{(1-\tau)^2}{C_1\tau}.
\end{equation}
Even with our certified exponent $1-\tau=2^{-40}$ and $C_1=14$, \eqref{eq:overview-quad-cap} caps $\kappa$ below $2^{-83.8}$; at the schedule $\eps=2^{-75}$, $c=2^{-56}$ of~\cite[(38)]{OAI26} the first term of~\eqref{eq:overview-quad-cap} is $2^{-181}$, giving $2^{-182}$, while the cap itself is ${\approx}\,2^{-104}$ at the printed structural constants. Figure~\ref{fig:pipeline} contrasts the architecture of~\cite{OAI26} with ours, highlighting the stages we remove and add.

\begin{figure}[t]
\setlength{\abovecaptionskip}{5pt}
\centering
\begin{tikzpicture}[
  font=\footnotesize,
  >=Latex,
  stage/.style={draw, rounded corners=2.5pt, minimum height=1.14cm, align=center, inner sep=3pt},
  oldbox/.style={stage, fill=red!5!white, draw=red!50!black, dashed},
  newbox/.style={stage, fill=blue!6!white, draw=blue!60!black, thick},
  io/.style={draw=black!65, rounded corners=2pt, fill=black!4, minimum height=0.96cm, align=center, inner sep=3pt},
  arr/.style={->, thick, draw=black!75}
]
  % Top row: Prior pipeline [OAI26]
  \node[anchor=west, font=\footnotesize\bfseries, text=red!55!black] at (0, 2.48)
    {Prior pipeline (\cite{HvdH21} architecture, networks of~\cite{OAI26}; $\kappa = 2^{-182}$, quadratically capped):};

  \node[io, text width=1.82cm] (in1) at (0.93, 1.28) {1-D input\\$S{=}s_1{\cdots}s_d$\\(odd primes)};
  \node[oldbox, text width=3.38cm] (old1) at (3.96, 1.28)
    {\textbf{CRT + resampling}\\primes $s_i \to 2^\ell$ grid\\$O(d^2)$ adjacent swaps};
  \node[oldbox, text width=3.38cm] (old2) at (7.78, 1.28)
    {\textbf{Chunk-swap layout $L$}\\$O(k/K)$ field swaps\\cost $O(V P K^{\tau-1})$};
  \node[oldbox, text width=3.48cm] (old3) at (11.65, 1.28)
    {\textbf{Bluestein + ring conv.}\\chirp in $\C[y]/(y^r{+}1)$\\$M$-amplified error};
  \node[io, text width=1.48cm] (out1) at (14.58, 1.28) {DFT of\\length $S$};

  \draw[arr] (in1) -- (old1);
  \draw[arr] (old1) -- (old2);
  \draw[arr] (old2) -- (old3);
  \draw[arr] (old3) -- (out1);

  % Shared foundation bar (single clean centered line <= 15.4cm)
  \node[draw=black!55, fill=yellow!10!white, rounded corners=2pt, minimum width=15.40cm, inner sep=3.2pt, align=center]
    (shared) at (7.70, 0.06)
    {\textbf{Shared $h{=}100$ network primitives~\cite{OAI26} (Fact~\ref{fact:oai-networks}):}\;\;
     interchange $O(V u^{\tau})$ \;\;$\bullet$\;\; butterfly layer $O(V d^{\lambda'})$};

  % Bottom row: This paper
  \node[anchor=west, font=\footnotesize\bfseries, text=blue!65!black] at (0, -0.68)
    {This paper (Theorems~\ref{thm:main}, \ref{thm:lean}; $\kappa = 2^{-46}$, linear in network deficit; $\kappa=2^{-56}$ with printed constants):};

  \node[io, text width=1.82cm] (in2) at (0.93, -1.90) {1-D input\\$R=2^k$\\(power of $2$)};
  \node[newbox, text width=3.38cm] (new1) at (3.96, -1.90)
    {\textbf{Fractional-shift bridge}\\half-bin kernel (\S\ref{sec:bridge})\\$O(d)$ swaps, same $2^k$ grid};
  \node[newbox, text width=3.38cm] (new2) at (7.78, -1.90)
    {\textbf{Walsh perm.\ (\S\ref{sec:walsh})}\\$L,(LB)^{-1}$ via $2^s H_T\,\CZ\,H_T$\\exact, cost $O(V\ell d^{\lambda'})$};
  \node[newbox, text width=3.48cm] (new3) at (11.65, -1.90)
    {\textbf{Direct DIF FFT (\S\ref{sec:direct-fft})}\\combined twiddles\\error $<12\ell\cdot 2^{-q}$};
  \node[io, text width=1.48cm] (out2) at (14.58, -1.90) {2-alias\\solve $\to F_R x$};

  \draw[arr] (in2) -- (new1);
  \draw[arr] (new1) -- (new2);
  \draw[arr] (new2) -- (new3);
  \draw[arr] (new3) -- (out2);
\end{tikzpicture}
\caption{Comparison of sub-$n\log n$ multiplication pipelines. Both build on the same two finite-network primitives of~\cite{OAI26}, but we replace the CRT/resampling front end, the chunk-swap address layout, and the Bluestein ring-convolution middle transform. Boxes removed by this paper: CRT map, prime lengths, Gaussian resampling, chunk-swap layout, Bluestein chirp, ring convolution; added: a fixed-kernel spread, an exact Walsh-layer permutation, and a two-term deconvolution.}
\label{fig:pipeline}
\end{figure}

\subsection{A CRT-free front end: what the bridge removes and how}\label{sec:overview-bridge}

\paragraph{Cooley--Tukey twiddles as fractional shifts.}
Consider first dimension two~\cite{MS}: factor $R=t_1 r$ into powers of two, write input indices as $j=j_1+t_1 j_2$ and reversed-digit frequencies as $f=u_1 r+u_2$ ($0\le j_1,u_1<t_1$, $0\le j_2,u_2<r$). Expanding the normalized transform $(F_R x)_f = \frac1R\sum_{j=0}^{R-1} x_j e^{2\pi i j f/R}$ yields
\begin{equation}\label{eq:overview-ct}
  (F_R x)_{u_1 r + u_2}
  = \frac{1}{t_1}\sum_{j_1=0}^{t_1-1} e^{2\pi i j_1 u_1/t_1}
    \Biggl(\frac{1}{r}\sum_{j_2=0}^{r-1} x_{j_1+t_1 j_2}\,e^{2\pi i (j_2 + j_1/t_1) u_2/r}\Biggr).
\end{equation}
The inner sum is the length-$r$ DFT of row~$j_1$, except that every sample is evaluated at the \emph{fractional shift} $j_2 + \delta_2$ with $\delta_2=j_1/t_1\in[0,1)$. Rather than multiplying by twiddles in frequency (which couples the axes), we spread each row $j_1$ in the time domain by a localized kernel $g(s)$ folded modulo~$r$ at offset $\delta_2$. By Poisson summation, the 2-D tensor DFT $F_{t_1,r}U$ of the spread array~$U$ is a weighted sum of aliases of $F_R x$ on the \emph{exact same} $t_1\times r$ grid:
\[
  (F_{t_1,r}U)_{u_1,u_2} = \sum_{\nu\in\Z} \hat g\Bigl(\frac{u_2}{r}-\nu\Bigr)\,(F_R x)_{((u_1-\nu)\bmod t_1)\,r + u_2}.
\]
For the sinc-Gaussian base kernel $g_0(s)=\sinc(s)e^{-\pi s^2/\alpha^2}$ ($\alpha\asymp\sqrt{\tw}$ for alias-width parameter~$\tw$), the Fourier transform $G(\xi)=\int_{\xi-1/2}^{\xi+1/2}\alpha e^{-\pi\alpha^2 v^2}\,dv$ is a smoothed indicator of $[-\frac12,\frac12]$. For $\xi\in(0,1)$, all aliases $\nu\notin\{0,1\}$ lie in Gaussian tails of total weight $O(2^{-\tw})$, leaving a two-term cyclic bidiagonal system $c_0 Z_{u_1} + c_1 Z_{u_1-1} = Y_{u_1}$ per column.

\paragraph{The Nyquist singularity and negative half-bin modulation.}
If the spreading kernel is real-valued, conjugate symmetry $\hat g(-\xi)=\overline{\hat g(\xi)}$ forces $|c_0|=|\hat g(\frac12)|=|\hat g(-\frac12)|=|c_1|$ at the Nyquist sample $u_2=r/2$ ($\xi=\frac12$). On a cycle of length~$L$, the bidiagonal matrix $c_0 I + c_1 T$ is then singular whenever $g$ is even and $L$ is even (since $c_0=c_1$), and for \emph{any} real Schwartz kernel its inverse has $\ell^\infty$ norm $\Omega(L)$ whenever invertible (Proposition~\ref{prop:real-kernel-fails}). In dimension two, $L=t_1$ and~\cite{MS} handles the single Nyquist column via a separate alternating-sum identity. In $d$ dimensions, however, cycle lengths reach $L=A_a=\prod_{i<a}t_i\le R/r$, so an $\Omega(L)$ conditioning loss on each of the $d-1$ spread axes would cost $\sum_{a=2}^d\log_2 A_a=\Theta(dk)$ extra bits of working precision. That overhead is superlinear in~$k$ once $d\to\infty$, destroying the $O(k)$-bit budget on which $V=\Theta(n)$ rests.

We eliminate the singularity uniformly across all dimensions and cycle lengths by modulating the kernel on each axis $a\in\{2,\dots,d\}$ with a \emph{negative half-bin phase}:
\begin{equation}\label{eq:overview-carrier}
  g_a(s) = e^{-\pi i s/t_a}\,\sinc(s)\,e^{-\pi s^2/\alpha^2}, \qquad \hat g_a(\xi) = G\Bigl(\xi + \frac{1}{2t_a}\Bigr).
\end{equation}
This shifts the frequency samples from integer bins $u_a/t_a$ to half-bin midpoints $y=(u_a+\frac12)/t_a\in(0,1)$ (Figure~\ref{fig:kernel-sampling}). Because $t_a\ge 2$ is even, no integer $u_a\in[t_a]$ ever lands on the crossing point $y=\frac12$ (which corresponds to $u_a=(t_a-1)/2\notin\Z$). Every sample satisfies $|y-\frac12|\ge \frac{1}{2t_a}$, ensuring a separation $|c_0-c_1|=|G(y)-G(y-1)|\ge \frac{7}{4t_a}$ and an inverse norm
\[
  \|B_a^{-1}\|_\infty \le \frac{1}{|c_0-c_1|} \le \frac47\,t_a \le t_a
\]
completely independent of the cycle length $L=A_a$ (Lemma~\ref{lem:bridge-B}). The negative sign in~\eqref{eq:overview-carrier} is essential: a positive carrier $e^{+\pi i s/t_a}$ shifts $u_a=0$ to $y=-\frac{1}{2t_a}\notin(0,1)$, encroaching on the omitted $\nu=-1$ passband and degrading the alias error from $O(2^{-\tw})$ to $\Omega(\tw^{-1/2}2^{-\tw/4})$ at $t_a=2$ (Proposition~\ref{prop:positive-fails}). Because $e^{-\pi i s/t_a}$ is anti-periodic of period $t_a$, periodizing $g_a$ folds $g_0$ with alternating signs $(-1)^{(e-s)/t_a}$.

\begin{figure}[t]
\setlength{\abovecaptionskip}{5pt}
\centering
\begin{tikzpicture}[xscale=8.0, yscale=2.95, font=\footnotesize, >=Latex]
  % Shaded regions outside [0,1] where samples approach omitted alias passbands
  \fill[black!6] (-0.23, 0) rectangle (0, 1.06);
  \fill[black!6] (1.0, 0) rectangle (1.23, 1.06);
  \node[font=\scriptsize, text=black!65, align=center] at (-0.118, 0.28) {$y<0$: nears\\$\nu=-1$ band};
  \node[font=\scriptsize, text=black!65, align=center] at (1.118, 0.28) {$y>1$: nears\\$\nu=+2$ band};

  % Axes
  \draw[->, thick, black!80] (-0.25, 0) -- (1.26, 0) node[right] {$y$};
  \draw[->, thick, black!80] (0, -0.04) -- (0, 1.26) node[above left, font=\scriptsize] {weight};
  \draw[dotted, black!40] (-0.22, 1.0) -- (1.22, 1.0);
  \draw[dotted, black!40] (-0.22, 0.5) -- (1.22, 0.5);
  \draw[black!80, thick] (-0.008, 1.0) -- (0.008, 1.0);
  \draw[black!80, thick] (-0.008, 0.5) -- (0.008, 0.5);

  % Exact curves c_0(y) = G(y) and c_1(y) = G(y-1) at \tw = 8 (\alpha\sqrt\pi = 2\sqrt{8\ln 2})
  \draw[very thick, blue!75!black] plot[smooth] coordinates {
    (-0.225,0.9665) (-0.200,0.9771) (-0.175,0.9848) (-0.150,0.9901) (-0.125,0.9937)
    (-0.100,0.9961) (-0.075,0.9976) (-0.050,0.9985) (-0.025,0.9990) (0.000,0.9991)
    (0.025,0.9990) (0.050,0.9985) (0.075,0.9976) (0.100,0.9961) (0.125,0.9937)
    (0.150,0.9901) (0.175,0.9848) (0.200,0.9771) (0.225,0.9665) (0.250,0.9521)
    (0.275,0.9330) (0.300,0.9086) (0.325,0.8781) (0.350,0.8411) (0.375,0.7975)
    (0.400,0.7473) (0.425,0.6913) (0.450,0.6304) (0.475,0.5661) (0.500,0.5000)
    (0.525,0.4339) (0.550,0.3696) (0.575,0.3087) (0.600,0.2527) (0.625,0.2025)
    (0.650,0.1589) (0.675,0.1219) (0.700,0.0914) (0.725,0.0670) (0.750,0.0479)
    (0.775,0.0335) (0.800,0.0229) (0.825,0.0152) (0.850,0.0099) (0.875,0.0063)
    (0.900,0.0039) (0.925,0.0023) (0.950,0.0014) (0.975,0.0008) (1.000,0.0004)
    (1.025,0.0002) (1.050,0.0001) (1.075,0.0001) (1.100,0.0000) (1.150,0.0000)
    (1.200,0.0000) (1.225,0.0000)
  };
  \draw[very thick, teal!75!black, dashed] plot[smooth] coordinates {
    (-0.225,0.0000) (-0.200,0.0000) (-0.150,0.0000) (-0.100,0.0000) (-0.075,0.0001)
    (-0.050,0.0001) (-0.025,0.0002) (0.000,0.0004) (0.025,0.0008) (0.050,0.0014)
    (0.075,0.0023) (0.100,0.0039) (0.125,0.0063) (0.150,0.0099) (0.175,0.0152)
    (0.200,0.0229) (0.225,0.0335) (0.250,0.0479) (0.275,0.0670) (0.300,0.0914)
    (0.325,0.1219) (0.350,0.1589) (0.375,0.2025) (0.400,0.2527) (0.425,0.3087)
    (0.450,0.3696) (0.475,0.4339) (0.500,0.5000) (0.525,0.5661) (0.550,0.6304)
    (0.575,0.6913) (0.600,0.7473) (0.625,0.7975) (0.650,0.8411) (0.675,0.8781)
    (0.700,0.9086) (0.725,0.9330) (0.750,0.9521) (0.775,0.9665) (0.800,0.9771)
    (0.825,0.9848) (0.850,0.9901) (0.875,0.9937) (0.900,0.9961) (0.925,0.9976)
    (0.950,0.9985) (0.975,0.9990) (1.000,0.9991) (1.025,0.9990) (1.050,0.9985)
    (1.075,0.9976) (1.100,0.9961) (1.125,0.9937) (1.150,0.9901) (1.175,0.9848)
    (1.200,0.9771) (1.225,0.9665)
  };

  % Nyquist crossing at y = 1/2 (integer bin u_a/4 for u_a=2)
  \draw[red!75!black, thick, densely dashed] (0.5, 0) -- (0.5, 1.12);
  \draw[red!75!black, very thick, fill=white] (0.5, 0.500) circle [x radius=0.011, y radius=0.026];
  \node[red!75!black, fill=white, inner sep=1pt, font=\scriptsize\bfseries, align=center] at (0.56, 1.19)
    {Nyquist crossing $y=\tfrac12$: $c_0=c_1$ (integer sample $u_a/t_a$ singular)};

  % y-axis labels placed AFTER banner, cleanly clear of curves and dotted lines
  \node[above left=1pt and 2.5pt, font=\scriptsize] at (0, 1.0) {$1$};
  \node[left=3pt, fill=black!6, inner sep=1pt, font=\scriptsize] at (-0.008, 0.5) {$\tfrac12$};

  % Four integer-bin ticks u_a/4 \in {0, 1/4, 1/2, 3/4} in red, plus black reference tick at y=1
  \foreach \ux/\ulab in {0/0, 0.25/\tfrac14, 0.5/\tfrac12, 0.75/\tfrac34} {
    \draw[red!65!black, thick] (\ux, -0.025) -- (\ux, 0.025);
    \node[below=3pt, font=\scriptsize, text=red!65!black] at (\ux, -0.01) {$\ulab$};
  }
  \draw[black!75, thick] (1.0, -0.025) -- (1.0, 0.025);
  \node[below=3pt, font=\scriptsize, text=black!75] at (1.0, -0.01) {$1$};

  % Half-bin samples y = (u_a + 1/2)/4 for t_a = 4, \tw = 8: 1/8, 3/8, 5/8, 7/8
  \foreach \hx/\czero/\cone/\ymax/\hlab in {
    0.125/0.9937/0.0063/0.9937/\tfrac18,
    0.375/0.7975/0.2025/0.7975/\tfrac38,
    0.625/0.2025/0.7975/0.7975/\tfrac58,
    0.875/0.0063/0.9937/0.9937/\tfrac78} {
    \draw[blue!45!black, dotted] (\hx, 0) -- (\hx, \ymax);
    \fill[blue!80!black] (\hx, \czero) circle [x radius=0.009, y radius=0.021];
    \fill[teal!80!black] (\hx, \cone) circle [x radius=0.009, y radius=0.021];
    \node[below=13pt, font=\scriptsize\bfseries, text=blue!75!black] at (\hx, 0) {$\hlab$};
  }

  % Curve labels placed AFTER dotted lines, strictly inside (0.125, 0.375) and (0.625, 0.875)
  \node[blue!75!black, fill=white, inner sep=1.2pt, font=\scriptsize\bfseries] at (0.245, 0.72) {$c_0(y)=G(y)$};
  \node[teal!75!black, fill=white, inner sep=1.2pt, font=\scriptsize\bfseries] at (0.755, 0.65) {$c_1(y)=G(y{-}1)$};

  % Separation arrows at y = 3/8 and y = 5/8 (|c_0 - c_1| = 0.5949 > 7/16 = 0.4375)
  \draw[<->, thick, blue!75!black] (0.375, 0.228) -- (0.375, 0.772);
  \draw[<->, thick, blue!75!black] (0.625, 0.228) -- (0.625, 0.772);
  \node[fill=white, inner sep=1.2pt, font=\scriptsize, text=blue!80!black] at (0.250, 0.47)
    {$|c_0{-}c_1|\ge\tfrac{7}{4t_a}$};
\end{tikzpicture}
\caption{Retained alias weights $c_0(y)=G(y)$ and $c_1(y)=G(y-1)$ plotted at $\tw=8$ for $t_a=4$. Integer-bin samples $u_a/t_a\in\{0,\tfrac14,\tfrac12,\tfrac34\}$ (red ticks) include the Nyquist crossing $y=\frac12$ (at $u_a=t_a/2$), where $c_0=c_1$, making the two-alias operator singular. Negative half-bin samples $y=(u_a+\frac12)/t_a\in\{\tfrac18,\tfrac38,\tfrac58,\tfrac78\}$ (solid dots) stay strictly inside $(0,1)$ and at distance $\ge\frac{1}{2t_a}$ from $\frac12$, giving $|c_0-c_1|\ge\frac{7}{4t_a}$ and uniform conditioning $\|B_a^{-1}\|_\infty\le \frac47 t_a$.}
\label{fig:kernel-sampling}
\end{figure}

\paragraph{Triangular non-commuting factorization in $d$ dimensions.}
For general $R=t_1\cdots t_d=2^k$, put $A_a=\prod_{i<a}t_i$, write time indices as $j=\sum_{i=1}^d j_i A_i$ ($j_1$ least significant) and frequencies in reversed digit order $f(\mathbf u)=\sum_{i=1}^d u_i R/A_{i+1}$ ($u_d$ least significant). On axis $a\ge 2$, the fractional shift $\delta_a(\mathbf j)=(\sum_{i<a}j_i A_i)/A_a\in[0,1)$ depends on the \emph{original} digits below~$a$. We therefore spread the axes in \emph{descending} order $a=d,d-1,\dots,2$, forming $U=\mathcal S_2\mathcal S_3\cdots\mathcal S_d x$, so that whenever $\mathcal S_a$ acts, coordinates $1,\dots,a-1$ are still untouched. Dual to this in frequency space, shifting $f$ by $-\nu R/A_a$ modifies only the more significant frequency digits $u_1,\dots,u_{a-1}$ while leaving $u_a,\dots,u_d$ fixed. Consequently, Poisson summation along all $d-1$ spread axes factors into an exact \emph{triangular} operator product (Lemma~\ref{lem:bridge-A}):
\begin{equation}\label{eq:overview-triangular}
  F_{\mathbf t}U = \cA_2\cA_3\cdots\cA_d(F_R x), \qquad
  (\cA_a Z)_f = \sum_{\nu\in\Z} G\Bigl(\frac{u_a(f)+1/2}{t_a}-\nu\Bigr)\,Z_{(f-\nu R/A_a)\bmod R}.
\end{equation}
The operators $\cA_a$ do not commute, so deconvolution applies the two-alias cyclic inverses $B_2^{-1},B_3^{-1},\dots,B_d^{-1}$ in ascending order. Because $\|(B_2\cdots B_d)^{-1}\|_\infty\le (4/7)^{d-1}(R/t_1)\le R/t_1$, working precision $P=\Theta(p+\log R)$ and grid precision $q=P+2$ control both alias and fixed-point errors to $O(2^{-p}\max(1,\|x\|_\infty))$ (Theorem~\ref{thm:bridge-precision}). Finally, exposing each axis~$a$ by a single direct field interchange with the innermost axis~$d$ rather than a chain of adjacent swaps reduces the exposure overhead from $O(d^2)$ to $O(d)$ interchanges of cost $O\bigl(V d(1+\ell^\tau)\bigr)$ (Lemma~\ref{lem:bridge-D}).

\subsection{Butterfly layers as exact positional permutations}\label{sec:overview-walsh}

The central new idea of our middle transform is that the simultaneous butterfly layer of Proposition~\ref{prop:prop14-param} can itself execute \emph{any} permutation of the bit positions of a binary record address, bit-exactly and at the cost of $O(1)$ butterfly layers per slice.

\paragraph{Selective Walsh layers via three-layer cancellation.}
Consider an array of $2^{m_0}$ disk-valued records on the $q$-grid ($m_0=\log_2(R/r)=O(q)$), arranged so that a \emph{slice} selects one bit position from each of $D\le d$ consecutive $K$-bit address chunks. One black-box call to Proposition~\ref{prop:prop14-param} applies $H_0=\frac12\bigl(\begin{smallmatrix}1&1\\1&-1\end{smallmatrix}\bigr)$ to \emph{all} $D$ slice positions simultaneously ($H_{\mathrm{all}}$); we cannot disable individual chunks inside its recursive complex-network schedule. Observe, however, two $2\times 2$ identities involving the phase gate $S=\diag(1,i)$:
\begin{equation}\label{eq:overview-h3}
  H_0^3 = \tfrac12\,H_0, \qquad (H_0 S)^3 = \tfrac{1+i}{4}\,I.
\end{equation}
Indeed, $H_0^2=\frac12 I$ immediately gives $H_0^3=\frac12 H_0$, while cubing $H_0 S = \frac12\bigl(\begin{smallmatrix}1&i\\1&-i\end{smallmatrix}\bigr)$ yields $(H_0 S)^2 = \frac14\bigl(\begin{smallmatrix}1+i & 1+i\\ 1-i & i-1\end{smallmatrix}\bigr)$ and $(H_0 S)^3 = \frac{1+i}{4}I$. Suppose we want to apply $H_0$ only on an arbitrary subset $T$ of the $D$ slice positions ($s=|T|$) and act as the identity on the complement $U$ ($q_{\mathrm{sel}}=D-s$). Before each of \emph{three} full layer calls $H_{\mathrm{all}}$, we scan the array and multiply every record at address~$x$ by the scalar phase $(S_U z)(x) = i^{\sum_{j\in U} x_j} z(x)$. Because operators on distinct address bits tensor and commute, $(H_{\mathrm{all}} S_U)^3$ acts as $\frac12 H_0$ on every bit in~$T$ and collapses to $\frac{1+i}{4}I$ on every bit in~$U$. Rescaling by $2^s[2(1-i)]^{q_{\mathrm{sel}}}$ recovers the selective operator $H_T$ exactly (Lemma~\ref{lem:walsh-select}; Figure~\ref{fig:walsh-gadget}(a)).

\paragraph{Batched controlled flips without address movement.}
Next, conjugating $\diag(1,-1)$ by $H_0$ gives $2 H_0\diag(1,-1)H_0 = \bigl(\begin{smallmatrix}0&1\\1&0\end{smallmatrix}\bigr)$, whereas $2 H_0 I H_0 = I$. Consider $s$ pairwise disjoint control-target bit pairs $(c_j,t_j)$ whose targets $T=\{t_1,\dots,t_s\}$ lie in one slice, while the controls $c_j$ may lie \emph{anywhere} in the $m_0$-bit record address. Define the recordwise sign diagonal $(\CZ\,z)(x)=(-1)^{\sum_{j=1}^s x_{c_j} x_{t_j}}z(x)$. Then (Lemma~\ref{lem:cnot}; Figure~\ref{fig:walsh-gadget}(b))
\begin{equation}\label{eq:overview-cnot}
  \CNOT^{\otimes s} = 2^s H_T\,\CZ\,H_T
\end{equation}
moves every record from address~$x$ to the address with $x_{t_j}\mapsto x_{t_j}\oplus x_{c_j}$ for all $j=1,\dots,s$, using six full butterfly-layer calls and linear scans (Proposition~\ref{prop:cnot-batch}). Crucially, $\CZ$ never routes data between $c_j$ and $t_j$: it computes the parity $\sum_j x_{c_j} x_{t_j}\bmod 2$ from each record's address in $q^{O(1)}$ steps (absorbed by the record payload of $\Theta(rq)$ bits, Lemma~\ref{lem:transfer-long-record}) and negates the record in place.

\begin{figure}[t]
\setlength{\abovecaptionskip}{5pt}
\centering
\begin{tikzpicture}[
  font=\footnotesize,
  >=Latex,
  gate/.style={draw=orange!70!black, fill=orange!15!yellow!20, thick, rounded corners=1.5pt, minimum width=0.50cm, minimum height=0.42cm, inner sep=1.5pt, align=center},
  layer/.style={draw=blue!70!black, fill=blue!10!white, thick, rounded corners=2pt, minimum width=0.72cm, inner sep=2pt, align=center},
  diagbox/.style={draw=teal!70!black, fill=teal!8!white, thick, rounded corners=2pt, inner sep=2.5pt, align=center},
  wire/.style={thick, black!75}
]
  % Panel (a): Selective Walsh gadget H_T
  \node[anchor=west, font=\footnotesize\bfseries] at (0, 2.55)
    {(a) Selective gadget $H_T$ ($s=|T|$, $q_{\mathrm{sel}}=|U|$):};

  % 4 wires: top 2 in T, bottom 2 in U
  \foreach \y/\wlab in {1.90/{$j\in T$}, 1.35/{$j'\in T$}, 0.65/{$u\in U$}, 0.10/{$u'\in U$}} {
    \draw[wire] (0.65, \y) -- (6.35, \y);
    \node[left, font=\scriptsize] at (0.65, \y) {\wlab};
  }

  % S gates on U wires before each of the 3 H_all layers
  \foreach \x in {1.10, 2.45, 3.80} {
    \node[gate] at (\x, 0.65) {$S$};
    \node[gate] at (\x, 0.10) {$S$};
  }
  % 3 full H_all boxes spanning all 4 wires
  \foreach \x in {1.78, 3.13, 4.48} {
    \node[layer, minimum height=2.25cm] at (\x, 1.00) {$H_{\mathrm{all}}$};
  }
  % Scalar box at end (wide enough so [2(1-i)]^{q_sel} fits cleanly)
  \node[diagbox, minimum height=2.25cm, minimum width=1.42cm, font=\scriptsize] at (5.62, 1.00)
    {$\times\,2^s\cdot$\\$[2(1{-}i)]^{q_{\mathrm{sel}}}$};

  \node[right, font=\scriptsize, align=left, text=blue!75!black] at (6.42, 1.62) {$H_0^3=\tfrac12 H_0$\\$\Rightarrow H_0$};
  \node[right, font=\scriptsize, align=left, text=orange!70!black] at (6.42, 0.38) {$(H_0 S)^3=\tfrac{1+i}{4}I$\\$\Rightarrow I$};

  % Panel (b): Batched CNOT = 2^s H_T CZ H_T
  \node[anchor=west, font=\footnotesize\bfseries] at (8.75, 2.55)
    {(b) Batched flips $\CNOT^{\otimes s}=2^s H_T\,\CZ\,H_T$:};

  % Control wires (top 2) and target wires in slice T (bottom 2)
  \draw[wire] (9.50, 1.90) -- (14.45, 1.90);
  \draw[wire] (9.50, 1.40) -- (14.45, 1.40);
  \draw[wire] (9.50, 0.65) -- (14.45, 0.65);
  \draw[wire] (9.50, 0.10) -- (14.45, 0.10);

  \node[left, font=\scriptsize] at (9.50, 1.90) {$x_{c_1}$};
  \node[left, font=\scriptsize] at (9.50, 1.40) {$x_{c_2}$};
  \node[left, font=\scriptsize] at (9.50, 0.65) {$x_{t_1}$};
  \node[left, font=\scriptsize] at (9.50, 0.10) {$x_{t_2}$};

  % First H_T on targets
  \node[layer, minimum height=1.05cm, minimum width=0.85cm] at (10.35, 0.375) {$H_T$};

  % CZ connections
  \draw[thick, teal!75!black] (11.40, 1.90) -- (11.40, 0.65);
  \fill[teal!75!black] (11.40, 1.90) circle (2.2pt);
  \fill[teal!75!black] (11.40, 0.65) circle (2.2pt);

  \draw[thick, teal!75!black] (11.90, 1.40) -- (11.90, 0.10);
  \fill[teal!75!black] (11.90, 1.40) circle (2.2pt);
  \fill[teal!75!black] (11.90, 0.10) circle (2.2pt);
  \node[above, font=\scriptsize, text=teal!75!black] at (11.65, 1.95) {$\CZ$};

  % Second H_T on targets + 2^s scaling
  \node[layer, minimum height=1.05cm, minimum width=0.85cm] at (12.85, 0.375) {$H_T$};
  \node[diagbox, minimum height=1.05cm, minimum width=0.78cm, font=\scriptsize] at (13.88, 0.375) {$\times\,2^s$};

  \node[right, font=\scriptsize] at (14.45, 1.90) {$x_{c_1}$};
  \node[right, font=\scriptsize] at (14.45, 1.40) {$x_{c_2}$};
  \node[right, font=\scriptsize, text=blue!75!black] at (14.45, 0.65) {$x_{t_1}\oplus x_{c_1}$};
  \node[right, font=\scriptsize, text=blue!75!black] at (14.45, 0.10) {$x_{t_2}\oplus x_{c_2}$};
\end{tikzpicture}
\caption{Turning simultaneous butterfly layers into exact controlled flips. (a)~Inserting $S=\diag(1,i)$ on $U$ before three full layers cancels $U$ to $\frac{1+i}{4}I$ while leaving $\frac12 H_0$ on $T$. (b)~Sandwiching the recordwise sign diagonal $\CZ$ between two selective $H_T$ gadgets flips all target address bits $t_j\in T$ by distant controls $c_j$ in six layer calls.}
\label{fig:walsh-gadget}
\end{figure}

\paragraph{Every positional permutation is $O(1)$ CNOT stages, and is bit-exact.}
Every permutation~$\pi$ of the $m_0$ record-address bit positions factors on its cycles as $\pi = I_1\circ I_0$, a product of two involutions, and each involution is a set of disjoint two-bit swaps implemented by three stages of disjoint CNOTs (Lemma~\ref{lem:walsh-involutions}). We partition the $m_0$ positions into $N_{\mathrm{slice}}\le m_0/d + K$ slices of $D\le d$ chunks, plus fewer than~$K$ leftover positions handled by linear two-digit stream swaps. Grouping each of the six CNOT stages by target slice then executes any positional permutation~$\pi$ in at most $36 N_{\mathrm{slice}}$ layer calls, costing $O\bigl(V[(m_0/d+K)d^{\lambda'}+K]\bigr)=O(V\ell d^{\lambda'})$ steps, the same order as the FFT butterfly rounds themselves (Theorem~\ref{thm:positional-perm}).

Why does fixed-point truncation introduce zero error across complex linear combinations? Before each CNOT batch on a $q$-grid disk array, we append $8d$ zero fractional bits to lift to precision $q_*=q+8d$ (Lemma~\ref{lem:walsh-lift}). Each normalized $D$-bit layer divides by $2^D$ ($D\le d$), while $S_U$ and $\CZ$ multiply by units in $\{\pm 1,\pm i\}$, so throughout both $H_T$ gadgets every intermediate entry lies on a dyadic grid of denominator exponent at most $q+s+3D\le q+4d < q_*$. Because Proposition~\ref{prop:prop14-param} evaluates each layer \emph{exactly} internally (using its guard~$\Delta$) before a single final truncation $Q_{q_*}$, every truncation discards only trailing zeros (Lemma~\ref{lem:walsh-exact}). Once $\CNOT^{\otimes s}$ completes, the array is an exact permutation of the original $q$-grid records, so stripping the $8d$ zero bits returns to the $q$-grid bit-exactly and resets denominator growth before the next batch.

\subsection{A direct tensor FFT and where \texorpdfstring{$2^{-46}$}{2\^{}\{-46\}} comes from}\label{sec:overview-fft}

\paragraph{Direct decimation-in-frequency tensor FFT.}
With exact positional permutations available at layer cost, Bluestein's chirp reduction and polynomial-ring convolutions over $\C[y]/(y^r+1)$ disappear completely. We evaluate $F_{\mathbf t}U$ on the $q$-grid ($t_d=r=2^\ell$, $t_i\in\{r/2,r\}$; Theorem~\ref{thm:direct-fft}) in four steps.
\begin{enumerate}[label=(\arabic*),nosep,leftmargin=*]
  \item Run a standard length-$r$ DIF FFT inside each contiguous record along axis~$d$.
  \item Apply the exact positional permutation~$L$ (Lemma~\ref{lem:walsh-layout}) to interleave $K$-bit chunks of the $d-1$ record axes.
  \item Execute the $\ell$ DIF bit-levels across all $d-1$ record axes in parallel: at most~$K$ prefix/leftover levels by exact individual butterflies in $O(VdK)$, and the remaining chunk levels by one call to Proposition~\ref{prop:prop14-param} per level in $O(V\ell d^{\lambda'})$. At each level~$v$, multiplying $d-1$ separately rounded twiddles would accumulate $O(d\ell\cdot 2^{-q})$ error. Instead, we combine all participating axis twiddles into a \emph{single} exact index $J(x)=\sum_i (r/2^{v+1})b_i j_i^{<v}\bmod r$, where $2^{v+1}$ is the DIF subproblem length, $b_i\in\{0,1\}$ is axis~$i$'s branch bit at level~$v$, and $j_i^{<v}$ is its lower $v$-bit index. Looking up that index in one $2r$-entry disk-valued phase table (Lemma~\ref{lem:phase-table}) incurs error $<6\cdot 2^{-q}$ per level regardless of~$d$.
  \item Apply the exact positional permutation $(LB)^{-1}=B^{-1}L^{-1}$ (Theorem~\ref{thm:positional-perm}) to undo both the chunk layout~$L$ and the DIF bit-reversal~$B$ in one pass.
\end{enumerate}
Because every exact DIF round is an $\ell^\infty$ contraction and the positional maps are exact, the $\ell$ last-axis levels and at most~$\ell$ record-axis rounds accumulate total error strictly below $12\ell\cdot 2^{-q}$ (Corollary~\ref{cor:walsh-middle}), with no $R/r$ error amplification and no output bit-widening.

\paragraph{Cost balance and the exponent $\kappa=2^{-46}$.}
Dividing the cost of one bridge call by $V=RP=\Theta(n)$ leaves the five normalized terms of Table~\ref{tab:cost-walsh}, where $P=\Theta(\log n)$, $d=\lfloor k^\eps\rfloor=\Theta(P^\eps)$, $\ell=\Theta(P^{1-\eps})$, $K=\lfloor d^c\rfloor=\Theta(P^{\eps c})$, and $\mathsf m(N)=\mathsf M(N)/N=(\log N)^{O(1)}$. Individual prefix butterfly rounds cost $O(dK)$ (gap $1-\eps(1+c)$ below $P^1$), simultaneous layers and the exact positional maps $L,(LB)^{-1}$ cost $O(\ell d^{\lambda'})$ (gap $\eps(1-\lambda')$), and the last-axis FFT and combined twiddle products cost $O(\ell\,\mathsf m(P))$ (gap $\eps$ up to $\log P$). Direct axis exposures and restorations cost $O\bigl(d(1+\ell^\tau)\bigr)$ (gap $(1-\eps)(1-\tau)$), and spreading and two-alias cyclic solves cost $O\bigl(d(\mathsf m(P)+\mathsf m(P^2))\bigr)$ (gap $1-\eps$ up to $\log P$). With the external chunk-swap term $P K^{\tau-1}$ gone, we are free to choose $c$ tiny so internal layer basis changes never bind. This leaves the butterfly-layer gap $g_{\mathrm{lay}}=\eps(1-\lambda') < (1-\sigma)/C_1$ as the sole bottleneck, now \emph{linear} in the complex network's deficit $1-\sigma$.

At $h=100$ ($m=10^6$), the complex network of~\cite{OAI26} has exact relative deficit $\eta_c=(W_c m - s_c)/(W_c m)=73/19{,}906{,}842{,}167{,}500$. Using the elementary rational bound $\ln 10^6 < 14$ (via $e^7>\sum_{j=0}^{11}7^j/j!>1000$, sharpening $\ln m<20$ in~\cite{OAI26}), Proposition~\ref{prop:deficits} certifies $s_c/W_c < m^\sigma$ at $\sigma=1-\xi$ with $\xi=2^{-42}$ (since $73\cdot 2^{42} > 14\cdot 19{,}906{,}842{,}167{,}500$) and $\tau=1-4\xi=1-2^{-40}$ for the bit network. Moreover, the guard bound $B_{\mathrm{net}}<10^{66}=m^{11}$ admits guard exponent $C_1=11+3=14$ (Proposition~\ref{prop:prop14-param}; cf.~\cite[\S5.4]{OAI26}). Setting $\eps=1/15$, $\beta=1/64$, $c=\xi/256$, $\lambda=1-\frac{63}{64}\xi$, and $\lambda'=1-\frac{31}{32}\xi$ gives $g_{\mathrm{lay}}=\eps(1-\lambda')=\frac{31}{480}\xi > \frac{1}{16}\xi = 2^{-46}$, yielding $\kappa=2^{-46}$ with positive slack (Section~\ref{sec:main-rows}). In short, Bluestein's chunk layout is capped quadratically at $(1-\tau)^2/(C_1\tau)\approx 2^{-83.8}$, whereas Walsh-layer permutations reach $(1-\sigma)/C_1\approx 2^{-45.8}$ linearly and already achieve $\kappa=2^{-56}$ under the printed constants of~\cite[Prop.~14]{OAI26}.

\subsection{Formal verification}\label{sec:overview-lean}

Section~\ref{sec:lean} describes our Lean~4 formalization~\cite{Lean26} of Theorem~\ref{thm:lean}. It constructs a single finite-alphabet multi-stack Turing machine in Mathlib's \texttt{Turing.TM2ComputableInTime} model and proves \texttt{SubNLogN.Challenge.IntMulSubNLogN} (exact multiplication of all $(a,b)\in\N^2$ within $C\,L\,(\log_2 L)^{1-2^{-46}}$ steps for input length $L\ge N_0$) from Lean's three standard axioms alone. That development verifies both $h=100$ finite networks, their rank and endpoint certificates, all stream and parameter-transfer procedures, and the analytical error bounds from scratch. The formal proposition does not express syntactic CRT-freeness and formalizes only the $\kappa=2^{-46}$ Walsh machine (deriving $2^{-47}$, $2^{-56}$, and $o(n\log n)$ by exponent monotonicity). It differs from the text primarily in two respects. First, it uses tie-upward nearest rounding (our bridge theorems are proved in Lean for arbitrary nearest rounding). Second, it instantiates the inner $O(N\log N)$ multiplier and its root offset via our independent Lean~4 formalization (\texttt{IntMulTM.\allowbreak intMul\_\allowbreak TM2\_\allowbreak nlogn}, kernel-checked within~\cite{Lean26}) of the two-dimensional multiplier of~\cite{MS}, which distributes its own separate Lean development.
